\documentclass[10pt]{article}

\usepackage[margin=0.75in]{geometry}
\usepackage{amsmath,amsthm,amssymb,bm}
\usepackage{xcolor}
\usepackage{cancel}
\usepackage{graphicx}
\usepackage{changepage}
\usepackage{circuitikz}
\usepackage{pgfplots}
\usepackage{physics}
\usepackage{hyperref}
\usepackage{siunitx}
\usepackage{minted}
\usepackage{fontspec}
\usepackage{relsize}
\usepackage{subfig}
\usepackage{mhchem}
\usepackage{todonotes}
\usepackage{biblatex}
\usepackage{multicol, multirow, booktabs}
\usepackage[breakable]{tcolorbox}
\usepackage[inline]{enumitem}
\patchcmd{\thebibliography}{\section*{\refname}}{}{}{}
\addbibresource{sources.bib}

\theoremstyle{definition}
\newtheorem{problem}{Problem}
\newtheorem{soln}{Solution}

\pgfplotsset{compat=newest}
\usetikzlibrary{lindenmayersystems}
\usetikzlibrary{arrows}
\usetikzlibrary{calc}
\usetikzlibrary{positioning, fit}
\usetikzlibrary{3d, perspective}

\definecolor{incolor}{HTML}{303F9F}
\definecolor{outcolor}{HTML}{D84315}
\definecolor{cellborder}{HTML}{CFCFCF}
\definecolor{cellbackground}{HTML}{F7F7F7}
\newcommand{\ui}{\hat{i}}
\newcommand{\uj}{\hat{j}}
\newcommand{\uk}{\hat{k}}
\newcommand{\ux}{\hat{x}}
\newcommand{\uy}{\hat{y}}
\newcommand{\uz}{\hat{z}}
\newcommand{\uv}[1]{\hat{\bm{{#1}}}}
\newcommand{\pr}[1]{{#1^\prime}}
\newcommand{\justif}[2]{&{#1}&\text{#2}}
\newcommand{\dul}[1]{\underline{\underline{#1}}}
\newcommand{\pdiff}[2][]{{\frac{\partial^{#1}}{\partial {{#2}^{#1}}}}}
% \newcommand{\dif}[2][]{{\frac{d^{#1}}{d{{#2}^{#1}}}}}
\pgfdeclarelayer{background}  
\pgfsetlayers{background,main}
\AtBeginDocument{\RenewCommandCopy\qty\SI}

\makeatletter
\newcommand{\boxspacing}{\kern\kvtcb@left@rule\kern\kvtcb@boxsep}
\makeatother
\newcommand{\prompt}[4]{
    \ttfamily\llap{{\color{#2}[#3]:\hspace{3pt}#4}}\vspace{-\baselineskip}
}

\newcommand{\thevenin}[2]{
  \begin{center}
    \begin{circuitikz} \draw
      (0,0) -- (2,0) to[battery1, l_=$V_{Th}\eq#1$] (2,2) 
      to[resistor, l_=$R_{Th}\eq#2$] (0,2)
      ;
      \draw [o-] (-.07,2.079);
      \draw [o-] (-.07,0.079);
    \end{circuitikz}
  \end{center}
}

\newcommand{\norton}[2]{
  \begin{center}
    \begin{circuitikz} \draw
      (0,0) -- (3,0) to[american current source, l_=$I_{N}\eq#1$] (3,2) -- (0,2) (2,0)
      to[resistor, l=$R_{N}\eq#2$] (2,2)
      ;
      \draw [o-] (-.07,2.079);
      \draw [o-] (-.07,0.079);
    \end{circuitikz}
  \end{center}
}

\newcommand{\highlight}[1]{\colorbox{yellow}{#1}}

\newcommand{\ti}[1]{\widetilde{#1}}

\newfontface{\Kaufmann}{Kaufmann}
\DeclareTextFontCommand{\kf}{\Kaufmann}
\newcommand{\scriptr}{\fontsize{12pt}{12pt}\kf{r}}

\newfontface{\KaufmannB}{Kaufmann Bold}
\DeclareTextFontCommand{\kfb}{\KaufmannB}
\newcommand{\bscriptr}{\fontsize{12pt}{12pt}\kfb{r}}

\newcommand{\bv}[1]{\vec{#1}}

\title{Physics 4605H: Assignment IX}
\author{Jeremy Favro (0805980) \\ Trent University, Peterborough, ON, Canada}
\date{\today}

\begin{document}
\maketitle

% PROBLEM 1
\begin{problem}
Practice with dual rail encoding:
Suppose that qubit 1 is encoded on rails 1 and 2, and qubit 2 is encoded on rails 3 and 4.
Now suppose that the state of qubits 1 and 2 is the 01 Bell state:
$$
  \ket{\psi}_{q_1q_2}=\frac{1}{\sqrt{2}}\left(\ket{00}_{q_1q_2}-\ket{11}_{q_1q_2}\right)
$$
Express this state in the rail basis, i.e. in terms of states $\ket{ijk\ell}_{1234}$ where $i, j, k, \ell = 0, 1$.
Ensure that in your solution every Dirac ket has a subscript specifying whether it is a qubit
state or a rail state.
\end{problem}
\begin{soln}
  Recall we defined $\ket{0}_q=\ket{01}_{12}=\hat{a}_2^\dagger\ket{00}_{12}$ as a photon in the lower rail (rail 2) and nothing in the upper rail (rail 1)
  and $\ket{1}_q=\ket{10}_{12}=\hat{a}_1^\dagger\ket{00}_{12}$ as the reverse. This means we can write the tensor product states as
  $$\ket{00}_{q_1q_2}=\ket{0}_{q_1}\otimes\ket{0}_{q_2}=\ket{0101}_{1234}=\hat{a}_2^\dagger\hat{a}4^\dagger\ket{0000}_{1234}$$
  and
  $$\ket{11}_{q_1q_2}=\ket{1}_{q_1}\otimes\ket{1}_{q_2}=\ket{1010}_{1234}=\hat{a}_1^\dagger\hat{a}_3^\dagger\ket{0000}_{1234}$$ so
  we can write the entire $\ket{\psi}_{q_1q_2}$ state as
  $$\ket{\psi}_{q_1q_2}=\frac{1}{\sqrt{2}}\left(\hat{a}_2^\dagger\hat{a}_4^\dagger-\hat{a}_1^\dagger\hat{a}_3^\dagger\right)\ket{0000}_{1234}.$$
\end{soln}
\newpage

% PROBLEM 2
\begin{problem}
Examining the effect of a phase shifter:
Regarding notation, here we will only be concerned with a single mode, so to simplify writing
let us replace $\hat{a}_{k,in}^\dagger$ with just $\hat{a}^\dagger$ and $\ket{0}_k$ with $\ket{0}$.
Meanwhile, $\hat{a}_{out}^\dagger\equiv \hat{U}\hat{a}^\dagger\hat{U}^\dagger\ket{0}$ describes the emitted state when the incident state was
$\hat{a}^\dagger\ket{0}$, where $\hat{U}=e^{-i\hat{H}t/\hbar}=e^{+i\phi\hat{a}^\dagger\hat{a}}$.
The goal here is to show that
$$\hat{a}^\dagger\ket{0}=e^{i\phi}\hat{a}^\dagger\ket{0}.$$
Here you will just consider the first few terms:
\begin{enumerate}[label=(\alph*)]
  \item Write $\hat{U}\hat{a}^\dagger\hat{U}^\dagger$ as $(\dots)\hat{a}^\dagger(\dots)$, writing out at least the first three terms of the Taylor expansion for
        each of the two exponentials.
  \item Re-express this as a sum of terms with different powers of $(i\phi)$, carefully writing the coefficient of $(i\phi)^0$, $(i\phi)^1$, and $(i\phi)^2$.
  \item Let $\hat{c}_n$ be the coefficient of $(i\phi)^n$. Use $\left[\hat{a},\hat{a}^\dagger\right]=1$ to show that $\hat{c}_n\ket{0}=A\hat{a}^\dagger\ket{0}$
        and find the scalar $A$ for $n=0,1,2$.
  \item Assuming that the pattern continues, complete the proof.
\end{enumerate}
\end{problem}
\begin{soln}
  \begin{enumerate}[label=(\alph*)]
    \item Recall that the Maclaurin series of $e^{\hat{A}}$ is
          $$e^{\hat{A}}=1+\hat{A}+\frac{1}{2}\hat{A}^2+\dots.$$
          The series we want then is, in the first three terms,
          $$e^{+i\phi\hat{a}^\dagger\hat{a}}\approx 1+i\phi\hat{a}^\dagger\hat{a}+\frac{1}{2}\left(i\phi\hat{a}^\dagger\hat{a}\right)^2.$$
          Hence
          \begin{align*}
            \hat{U}\hat{a}^\dagger\hat{U}^\dagger & \approx \left[1+i\phi\hat{a}^\dagger\hat{a}+\frac{1}{2}\left(i\phi\hat{a}^\dagger\hat{a}\right)^2\right]\hat{a}^\dagger\left[1-i\phi\hat{a}^\dagger\hat{a}+\frac{1}{2}\left(-i\phi\hat{a}^\dagger\hat{a}\right)^2\right]
          \end{align*}
    \item Expanding the above expression,
          \begin{align*}
            \hat{U}\hat{a}^\dagger\hat{U}^\dagger & \approx \left[1+i\phi\hat{a}^\dagger\hat{a}+\frac{1}{2}\left(i\phi\hat{a}^\dagger\hat{a}\right)^2\right]\hat{a}^\dagger\left[1-i\phi\hat{a}^\dagger\hat{a}+\frac{1}{2}\left(-i\phi\hat{a}^\dagger\hat{a}\right)^2\right]                                             \\
                                                  & = \left[1+i\phi\hat{a}^\dagger\hat{a}+\frac{1}{2}(i\phi)^2\hat{a}^\dagger\hat{a}\hat{a}^\dagger\hat{a}\right]\left[\hat{a}^\dagger-i\phi\hat{a}^\dagger\hat{a}^\dagger\hat{a}+\frac{1}{2}(i\phi)^2\hat{a}^\dagger\hat{a}^\dagger\hat{a}\hat{a}^\dagger\hat{a}\right] \\
                                                  & = \left[\hat{a}^\dagger-i\phi\hat{a}^\dagger\hat{a}^\dagger\hat{a}+\frac{1}{2}(i\phi)^2\hat{a}^\dagger\hat{a}^\dagger\hat{a}\hat{a}^\dagger\hat{a}\right]                                                                                                            \\
                                                  & \quad +i\phi\hat{a}^\dagger\hat{a}\left[\hat{a}^\dagger-i\phi\hat{a}^\dagger\hat{a}^\dagger\hat{a}+\frac{1}{2}(i\phi)^2\hat{a}^\dagger\hat{a}^\dagger\hat{a}\hat{a}^\dagger\hat{a}\right]                                                                            \\
                                                  & \quad +\frac{1}{2}(i\phi)^2\hat{a}^\dagger\hat{a}\hat{a}^\dagger\hat{a}\left[\hat{a}^\dagger-i\phi\hat{a}^\dagger\hat{a}^\dagger\hat{a}+\frac{1}{2}(i\phi)^2\hat{a}^\dagger\hat{a}^\dagger\hat{a}\hat{a}^\dagger\hat{a}\right]                                       \\
                                                  & =
            \left[
              \hat{a}^\dagger
              -i\phi\hat{a}^\dagger\hat{a}^\dagger\hat{a}
              +\frac{1}{2}(i\phi)^2\hat{a}^\dagger\hat{a}^\dagger\hat{a}\hat{a}^\dagger\hat{a}
            \right]                                                                                                                                                                                                                                                                                                      \\
                                                  & \quad +
            \left[
              i\phi\hat{a}^\dagger\hat{a}\hat{a}^\dagger
              -(i\phi)^2\hat{a}^\dagger\hat{a}\hat{a}^\dagger\hat{a}^\dagger\hat{a}
              +\frac{1}{2}(i\phi)^3\hat{a}^\dagger\hat{a}\hat{a}^\dagger\hat{a}^\dagger\hat{a}\hat{a}^\dagger\hat{a}
            \right]                                                                                                                                                                                                                                                                                                      \\
                                                  & \quad +
            \left[
              \frac{1}{2}(i\phi)^2\hat{a}^\dagger\hat{a}\hat{a}^\dagger\hat{a}\hat{a}^\dagger
              -\frac{1}{2}(i\phi)^3\hat{a}^\dagger\hat{a}\hat{a}^\dagger\hat{a}\hat{a}^\dagger\hat{a}^\dagger\hat{a}
              +\frac{1}{4}(i\phi)^4\hat{a}^\dagger\hat{a}\hat{a}^\dagger\hat{a}\hat{a}^\dagger\hat{a}^\dagger\hat{a}\hat{a}^\dagger\hat{a}
            \right]                                                                                                                                                                                                                                                                                                      \\
                                                  & =
            \hat{a}^\dagger
            +i\phi\hat{a}^\dagger(\hat{a}\hat{a}^\dagger-\hat{a}^\dagger\hat{a})
            +(i\phi)^2\hat{a}^\dagger\left(\frac{1}{2}\hat{a}^\dagger\hat{a}\hat{a}^\dagger\hat{a}
            -\hat{a}\hat{a}^\dagger\hat{a}^\dagger\hat{a}
            +\frac{1}{2}\hat{a}\hat{a}^\dagger\hat{a}\hat{a}^\dagger
            \right)                                                                                                                                                                                                                                                                                                      \\
                                                  & \quad +\frac{1}{2}(i\phi)^3\left(
            \hat{a}^\dagger\hat{a}\hat{a}^\dagger\hat{a}^\dagger\hat{a}\hat{a}^\dagger\hat{a}
            -(i\phi)^3\hat{a}^\dagger\hat{a}\hat{a}^\dagger\hat{a}\hat{a}^\dagger\hat{a}^\dagger\hat{a}
            \right)                                                                                                                                                                                                                                                                                                      \\
                                                  & \quad +\frac{1}{4}(i\phi)^4\hat{a}^\dagger\hat{a}\hat{a}^\dagger\hat{a}\hat{a}^\dagger\hat{a}^\dagger\hat{a}\hat{a}^\dagger\hat{a}                                                                                                                                   \\
          \end{align*}
          Now $(\hat{a}\hat{a}^\dagger-\hat{a}^\dagger\hat{a})=\left[\hat{a},\hat{a}^\dagger\right]=1$ but what about
          $$\frac{1}{2}\hat{a}^\dagger\hat{a}\hat{a}^\dagger\hat{a}
            -\hat{a}\hat{a}^\dagger\hat{a}^\dagger\hat{a}
            +\frac{1}{2}\hat{a}\hat{a}^\dagger\hat{a}\hat{a}^\dagger
            ?$$
          Well we can rewrite it as
          \begin{align*}
             & =\frac{1}{2}\hat{a}^\dagger\hat{a}\hat{a}^\dagger\hat{a}
            -\hat{a}\hat{a}^\dagger\hat{a}^\dagger\hat{a}
            +\frac{1}{2}\hat{a}\hat{a}^\dagger\hat{a}\hat{a}^\dagger                                                            \\
             & =\frac{1}{2}\hat{a}^\dagger\hat{a}\hat{a}^\dagger\hat{a}
            +\hat{a}\hat{a}^\dagger\left(\frac{1}{2}\hat{a}\hat{a}^\dagger-\hat{a}^\dagger\hat{a}\right)                        \\
             & =\frac{1}{2}\hat{a}^\dagger\hat{a}\hat{a}^\dagger\hat{a}
            +\frac{1}{2}\hat{a}\hat{a}^\dagger\left(\hat{a}\hat{a}^\dagger-\hat{a}^\dagger\hat{a}-\hat{a}^\dagger\hat{a}\right) \\
             & =\frac{1}{2}\hat{a}^\dagger\hat{a}\hat{a}^\dagger\hat{a}
            +\frac{1}{2}\hat{a}\hat{a}^\dagger\left(1-\hat{a}^\dagger\hat{a}\right)                                             \\
             & =\frac{1}{2}\hat{a}^\dagger\hat{a}\hat{a}^\dagger\hat{a}-\frac{1}{2}\hat{a}\hat{a}^\dagger\hat{a}^\dagger\hat{a}
            +\frac{1}{2}\hat{a}\hat{a}^\dagger                                                                                  \\
             & =\frac{1}{2}\left(\hat{a}^\dagger\hat{a}-\hat{a}\hat{a}^\dagger\right)\hat{a}^\dagger\hat{a}
            +\frac{1}{2}\hat{a}\hat{a}^\dagger                                                                                  \\
             & =-\frac{1}{2}\left(\hat{a}\hat{a}^\dagger-\hat{a}^\dagger\hat{a}\right)\hat{a}^\dagger\hat{a}
            +\frac{1}{2}\hat{a}\hat{a}^\dagger                                                                                  \\
             & =-\frac{1}{2}\hat{a}^\dagger\hat{a}
            +\frac{1}{2}\hat{a}\hat{a}^\dagger                                                                                  \\
             & =\frac{1}{2}
          \end{align*}
          where I think the algebra is clear at every step. We're just trying to massage the expression to get terms that look like $\hat{a}\hat{a}^\dagger-\hat{a}^\dagger\hat{a}=\left[\hat{a},\hat{a}^\dagger\right]=1$.
          I've also just realized I'm basically doing part (c) here.
    \item See above. The scalars are $1$, $i\phi$, $(i\phi)^2/2$.
    \item Assuming the pattern continues this is just the Taylor series
          $$e^{i\phi}=1+i\phi+\frac{1}{2}(i\phi)^2$$
          multiplied by $\hat{a}^\dagger$, as we wanted.
  \end{enumerate}
\end{soln}

% PROBLEM 3
\begin{problem}
Suppose that Alice and Bob implement quantum teleportation using qubits
$1$, $2$, and $3$, and Alice begins with qubit $1$ in state $\alpha\ket{0}_1+\beta\ket{1}_1$ while qubits $2$ and $3$ are in
the $11$ Bell state: $(\ket{01}_{23}-\ket{10}_{23})/\sqrt{2}$. (A full list of Bell states appear on pp.41-42 of the
text.)
\begin{enumerate}[label=(\alph*)]
  \item Alice performs her Bell measurement on qubits $1$ and $2$, determines they are in Bell state
        $\ket{B_{nm}}_{12}$ ($n, m = 0, 1$) and sends the values of $n$ and $m$ to Bob. For each possible $nm$ pair,
        determine what gate Bob must apply to qubit 3 to complete the teleportation of Alice's
        qubit.
  \item What is the final state of Alice's qubit?
\end{enumerate}
\end{problem}
\begin{soln}~
  \begin{enumerate}[label=(\alph*)]
    \item I want to try to actually do the math a bit. I'm not sure if there is an easier way I'm forgetting or what.
          We'll relabel the qubit to be teleported as $\ket{\psi}_{T}=\alpha\ket{0}_T+\beta\ket{1}_T$ and the qubits of Alice and Bob respectively as
          $\ket{\psi}_A$ and $\ket{\psi}_B$. We are saying that the $A$ and $B$ qubits are in Bell state
          $$\ket{\Psi}_{AB}=\frac{1}{\sqrt{2}}\left(\ket{0}_A\otimes\ket{1}_B-\ket{1}_A\otimes \ket{0}_B\right).$$
          This means that the state of the entire three qubit system is given by
          $$\ket{\psi}_{T}\otimes \ket{\Psi}_{AB}=\frac{1}{\sqrt{2}}\left(\alpha\ket{0}_T+\beta\ket{1}_T\right)\otimes\left(\ket{0}_A\otimes\ket{1}_B-\ket{1}_A\otimes \ket{0}_B\right).$$
          Alice then makes her measurement in the Bell basis on qubits $\ket{\psi}_A$ and $\ket{\psi}_{T}$. We can expand the above expression as 
          $$\ket{\psi}_{T}\otimes \ket{\Psi}_{AB}=
          \frac{1}{\sqrt{2}}\left(\alpha\ket{0}_T\otimes\ket{0}_A\otimes\ket{1}_B-\alpha\ket{0}_T\otimes\ket{1}_A\otimes \ket{0}_B
          +\beta\ket{1}_T\otimes\ket{0}_A\otimes\ket{1}_B-\beta\ket{1}_T\otimes\ket{1}_A\otimes \ket{0}_B\right).$$
          To get this into a form that's more useful we can rewrite the tensor products of $T$ and $A$ qubits by rearranging the Bell state expressions, for example
          $$\ket{0}_T\otimes\ket{0}_A=\frac{1}{\sqrt{2}}\left(\ket{B_{00}}+\ket{B_{01}}\right)$$
          using the Bell state notation from the text. This lets us write 
          \begin{align*}
            \ket{\psi}_{T}\otimes \ket{\Psi}_{AB}&=
          \frac{1}{\sqrt{2}}\left(\frac{\alpha}{\sqrt{2}}\left(\ket{B_{00}}_{TA}+\ket{B_{01}}_{TA}\right)\otimes\ket{1}_B\right.\\
          &\quad-\frac{\alpha}{\sqrt{2}}\left(\ket{B_{10}}_{TA}+\ket{B_{11}}_{TA}\right)\otimes \ket{0}_B\\
          &\quad+\frac{\beta}{\sqrt{2}}\left(\ket{B_{10}}_{TA}-\ket{B_{11}}_{TA}\right)\otimes\ket{1}_B\\
          &\quad\left.-\frac{\beta}{\sqrt{2}}\left(\ket{B_{00}}_{TA}-\ket{B_{01}}_{TA}\right)\otimes \ket{0}_B\right)\\
          &=
          \frac{1}{2}\left(\ket{B_{00}}_{TA}\otimes\alpha\ket{1}_B+\ket{B_{01}}_{TA}\otimes\alpha\ket{1}_B\right.\\
          &\quad -\ket{B_{10}}_{TA}\otimes \alpha\ket{0}_B-\ket{B_{11}}_{TA}\otimes \alpha\ket{0}_B\\
          &\quad+\ket{B_{10}}_{TA}\otimes\beta\ket{1}_B-\ket{B_{11}}_{TA}\otimes\beta\ket{1}_B\\
          &\quad\left.-\ket{B_{00}}_{TA}\otimes \beta\ket{0}_B+\ket{B_{01}}_{TA}\otimes \beta\ket{0}_B\right)\\
          &=
          \frac{1}{2}\left(\ket{B_{00}}_{TA}\otimes\left(\alpha\ket{1}_B-\beta\ket{0}_B\right)+\ket{B_{01}}_{TA}\otimes\left(\alpha\ket{1}_B+\beta\ket{0}_B\right)\right.\\
          &\quad \ket{B_{10}}_{TA}\otimes \left(-\alpha\ket{0}_B+\beta\ket{1}_B\right)-\ket{B_{11}}_{TA}\otimes \left(\alpha\ket{0}_B+\beta\ket{1}_B\right).
          \end{align*}
          Alice can measure one of the $\ket{B_{nm}}_{TA}$ states, each of which corresponds to some variation on the state we want to teleport, so we need to find the gates to
          correct this.
          \begin{enumerate}[label=\roman*.]
            \item If Alice measures $\ket{B_{00}}_{TA}$ then we are in $\alpha\ket{1}_B-\beta\ket{0}_B$ and have to apply $\hat{\sigma}_x$, then 
            $\hat{\sigma}_y$. The sequence $\hat{\sigma}_y\hat{\sigma}_x$ is equivalent to $-i\hat{\sigma}_z$.
            \item If Alice measures $\ket{B_{01}}_{TA}$ then we are in $\alpha\ket{1}_B+\beta\ket{0}_B$ and we just have to apply $\hat{\sigma}_x$.
            \item If Alice measures $\ket{B_{10}}_{TA}$ then we are in $-\alpha\ket{0}_B+\beta\ket{1}_B$ and need to apply $\hat{\sigma}_x$ to get 
            $-\alpha\ket{1}_B+\beta\ket{0}_B$, then $\hat{\sigma}_z$ to get $\alpha\ket{1}_B+\beta\ket{0}_B$, then $\hat{\sigma}_x$ again to get to 
            $\alpha\ket{0}_B+\beta\ket{1}_B$. The sequence $\hat{\sigma}_x\hat{\sigma}_z\hat{\sigma}_x$ is equivalent to $-i\hat{\sigma}_y\hat{\sigma}_x$ which is itself
            equivalent to $-\hat{\sigma}_z$.
            \item If Alice measures $\ket{B_{11}}_{TA}$ then we are in $-\alpha\ket{0}_B-\beta\ket{1}_B$ and must apply $\hat{\sigma}_z$ to obtain $-\alpha\ket{0}_B+\beta\ket{1}_B$,
            then $\hat{\sigma}_x$ to obtain $-\alpha\ket{1}_B+\beta\ket{0}_B$, then $\hat{\sigma}_z$ again, and finally $\hat{\sigma}_x$ to obtain 
            $\alpha\ket{0}_B+\beta\ket{1}_B$. The sequence $\hat{\sigma}_x\hat{\sigma}_z\hat{\sigma}_x\hat{\sigma}_z=-\sigma_y^2=-\mathbb{I}$.
          \end{enumerate}
    \item I think it would be an equal superposition between $\ket{0}$ and $\ket{1}$, because all she's measured it to be is in a specific Bell state with the $T$ qubit.
  \end{enumerate}
\end{soln}
\end{document}